% AUTO-GENERATED. Do not edit; generated from frozen science_contract.json.
\subsection*{Machine-readable science contract}
\begin{description}
\item[Research question] For source-native E010/PKLSA-qualified vortex carriers evolved by a preregistered finite-core filament or incompressible Euler provider, do the extracted paired Kelvin residual phase and mode envelope require Sine-Gordon restoring dynamics and/or Mittag-Leffler fractional-memory relaxation?
\item[Objective] Close the v0.2.3 provider gap by generating phi(s,t) and R(t) directly from qualified E010 carriers under frozen finite-core/Euler dynamics, then apply the canonical A056 blind competition without promoting simulation evidence to physical cross-source support.
\item[$H_0$] H0: Sine-Gordon and Mittag-Leffler descriptions do not achieve the preregistered discovery, held-out, numerical-certification and independent-replication thresholds beyond the registered alternatives.
\item[$H_1$] H1: one or both candidate mechanisms outperform the registered alternatives, survive held-out and numerical certification, and replicate across eligible independent physical evidence.
\end{description}
\subsubsection*{Registered equations}
\paragraph{E1: Wave null phase model}
\begin{equation}
\varphi_{tt}=a\varphi_{ss}
\end{equation}
\noindent\textit{Dimensional check:} Both sides have dimensions T\textasciicircum{}\{-2\} for dimensionless phase when [a]=L\textasciicircum{}2 T\textasciicircum{}\{-2\}.\\
\paragraph{E2: Klein-Gordon competitor}
\begin{equation}
\varphi_{tt}=a\varphi_{ss}-b\varphi
\end{equation}
\noindent\textit{Dimensional check:} [b]=T\textasciicircum{}\{-2\}; every term has dimensions T\textasciicircum{}\{-2\}.\\
\paragraph{E3: Duffing competitor}
\begin{equation}
\varphi_{tt}=a\varphi_{ss}-b\varphi-c\varphi^3
\end{equation}
\noindent\textit{Dimensional check:} For dimensionless phase, [b]=[c]=T\textasciicircum{}\{-2\}.\\
\paragraph{E4: Sine-Gordon candidate}
\begin{equation}
\varphi_{tt}=a\varphi_{ss}-b\sin\varphi
\end{equation}
\noindent\textit{Dimensional check:} sin(varphi) is dimensionless and [b]=T\textasciicircum{}\{-2\}.\\
\paragraph{E5: Bayesian information criterion used on discovery blocks}
\begin{equation}
\mathrm{BIC}=n\ln(\mathrm{RSS}/n)+k\ln n
\end{equation}
\noindent\textit{Dimensional check:} Model residuals are compared within one observable/scaling; BIC differences are dimensionless.\\
\paragraph{E6: Held-out normalized root-mean-square error}
\begin{equation}
\mathrm{NRMSE}=\sqrt{n^{-1}\sum_i(y_i-\hat y_i)^2}/\sigma_y
\end{equation}
\noindent\textit{Dimensional check:} Numerator and denominator have the same observable units, so NRMSE is dimensionless.\\
\paragraph{E7: Mittag-Leffler function}
\begin{equation}
E_{\alpha,\beta}(z)=\sum_{k=0}^{\infty} z^k/\Gamma(\alpha k+\beta)
\end{equation}
\noindent\textit{Dimensional check:} z must be dimensionless.\\
\paragraph{E8: Mittag-Leffler relaxation candidate}
\begin{equation}
R_{ML}(t)=E_{\alpha,1}[-(t/\tau)^\alpha]
\end{equation}
\noindent\textit{Dimensional check:} t/tau is dimensionless.\\
\paragraph{E9: Ordinary exponential null limit}
\begin{equation}
R_{exp}(t)=\exp(-t/\tau)=E_{1,1}(-t/\tau)
\end{equation}
\noindent\textit{Dimensional check:} t/tau is dimensionless.\\
\paragraph{E10: Stretched-exponential competitor}
\begin{equation}
R_{SE}(t)=\exp[-(t/\tau)^\beta]
\end{equation}
\noindent\textit{Dimensional check:} t/tau is dimensionless.\\
\paragraph{E11: Bi-exponential competitor}
\begin{equation}
R_{2e}(t)=w e^{-t/\tau_1}+(1-w)e^{-t/\tau_2}
\end{equation}
\noindent\textit{Dimensional check:} Each exponential argument is dimensionless and w is dimensionless.\\
\paragraph{E12: Phase-resolution convergence diagnostic}
\begin{equation}
\delta_c=\lVert\mathbf c_f-\mathbf c_c\rVert_2/\max(\lVert\mathbf c_f\rVert_2,10^{-30}),\quad \mathbf c=(a,b)
\end{equation}
\noindent\textit{Dimensional check:} Coefficient vectors are compared componentwise under the same coordinate units; ratio is dimensionless.\\
\paragraph{E13: Memory-resolution convergence diagnostics}
\begin{equation}
\delta_\alpha=|\alpha_f-\alpha_c|,\qquad \delta_\tau=|\tau_f-\tau_c|/\max(|\tau_f|,10^{-30})
\end{equation}
\noindent\textit{Dimensional check:} Both diagnostics are dimensionless.\\
\paragraph{E14: Incompressible unforced Euler evolution}
\begin{equation}
\partial_t\mathbf u+(\mathbf u\cdot\nabla)\mathbf u=-\nabla p,\qquad \nabla\cdot\mathbf u=0
\end{equation}
\noindent\textit{Dimensional check:} Each acceleration term has dimensions L T\textasciicircum{}\{-2\}; incompressibility has T\textasciicircum{}\{-1\}.\\
\paragraph{E15: Finite-core centerline vorticity seed}
\begin{equation}
\boldsymbol\omega(\mathbf x,0)\propto\sum_j \mathbf T_j\exp[-d_{per}(\mathbf x,\mathbf X_j)^2/(2\sigma^2)]
\end{equation}
\noindent\textit{Dimensional check:} The exponential is dimensionless; normalization fixes the velocity scale.\\
\paragraph{E16: Material-marker transport}
\begin{equation}
\dot{\mathbf X}(s,t)=\mathbf u(\mathbf X(s,t),t)
\end{equation}
\noindent\textit{Dimensional check:} Both sides have dimensions L T\textasciicircum{}\{-1\}.\\
\paragraph{E17: Frozen Kelvin perturbation}
\begin{equation}
\mathbf X_\epsilon(s,0)=\mathbf X_0(s)+\epsilon[\cos(m\theta)\mathbf e_1+\sin(m\theta)\mathbf e_2]
\end{equation}
\noindent\textit{Dimensional check:} epsilon has dimensions L; trigonometric factors and frame vectors are dimensionless.\\
\paragraph{E18: Complex paired transverse displacement}
\begin{equation}
q(s,t)=\delta\mathbf X\cdot\mathbf e_1+i\,\delta\mathbf X\cdot\mathbf e_2
\end{equation}
\noindent\textit{Dimensional check:} q has dimensions L.\\
\paragraph{E19: Residual Kelvin phase}
\begin{equation}
\varphi(s,t)=\arg[q(s,t)e^{-im\theta(s)}]
\end{equation}
\noindent\textit{Dimensional check:} The argument is an angle and therefore dimensionless.\\
\paragraph{E20: Projected mode envelope}
\begin{equation}
R(t)=|c_m(t)|/|c_m(0)|,\qquad c_m=N_s^{-1}\sum_j q(s_j,t)e^{-im\theta_j}
\end{equation}
\noindent\textit{Dimensional check:} R is dimensionless.\\
\subsubsection*{Registered code-to-mathematics steps}
\begin{enumerate}
\item \textbf{S1} [G0; equations E1]: Verify frozen protocol/blind commitments and enumerate hashed provider cases.
\item \textbf{S0P} [G0; equations E14, E15, E16, E17, E18, E19, E20]: Resolve and re-hash E010 v0.3.1 strict upstream-provider representatives, construct finite-core base/perturbed carriers, evolve them with the frozen provider, and commit the private carrier reveal before A056 scoring.
\item \textbf{S2} [G1; equations E1]: Validate provider metadata, array shapes, monotonicity, finite values and sampling uniformity.
\item \textbf{S3} [G2; equations E4, E9]: Evaluate Python FP64 finite differences, POD spectral stability and the Mittag-Leffler alpha=1 exponential identity.
\item \textbf{S4} [G3; equations E1, E2, E3, E4, E5, E8, E9, E10, E11]: Fit Wave/Klein-Gordon/Duffing/Sine-Gordon and exponential/stretched/bi-exponential/Mittag-Leffler families on discovery blocks and evaluate BIC.
\item \textbf{S5} [G4; equations E6]: Evaluate frozen held-out temporal blocks with NRMSE ratios.
\item \textbf{S6} [G5; equations E12, E13]: Apply deterministic coarsening diagnostics and strict C++/OpenMP FP64 parity to confirmed candidates.
\item \textbf{S7} [G6; equations E12]: Consume optional canonical-framework SYCL/DD32 smoke evidence as screening-only diagnostics.
\item \textbf{S8} [G7; equations E5, E6]: Assess evidence-class-aware replication; synthetic/simulation controls cannot close physical cross-source replication.
\item \textbf{S9} [G8; equations E4, E8]: Classify joint, phase-only or memory-only physical support only after eligible replication.
\end{enumerate}
